\section{Síntesis y ampliación}

Esta sección condensa el EP102 en siete conclusiones. Primera: el momento GPT de la CV no puede copiar sin más el de la NLP, porque las imágenes estáticas no traen de forma natural semántica humana ni alineación. Segunda: la dificultad central de la integración multimodal es que la generación y la comprensión no pueden corregirse entre sí. Tercera: la next token prediction es la base del despegue de los grandes modelos y también una de las causas de los saltos de pasos en el razonamiento y de la falta de precisión. Cuarta: la aportación principal de o1 es usar RL para activar la CoT, la reflexión y la búsqueda de critical decisions. Quinta: la Long CoT visual traslada el pensamiento lento al espacio visual. Sexta: el long context no es solo un problema de complejidad de la arquitectura, sino sobre todo de compresión, recuperación, aislamiento de escenarios y colaboración entre varios modelos. Séptima: el siguiente momento Agent, más profundo, podría venir del aprendizaje en línea y del environmental scaling.

\begin{importantbox}{El EP102 dentro de la serie de IA e internet de Zhang Xiaojun (张小珺)}
El EP102 es una de las entrevistas sobre la ruta multimodal con mayor contenido técnico de este lote. Se conecta con el informe técnico sobre Agents del EP110, con el agentic LLM del EP113 y con la filosofía de investigación sobre Agents del EP115; además, ofrece para los problemas de inteligencia corporizada de los EP106 y EP109 una explicación previa sobre razonamiento visual y modelos del mundo.
\end{importantbox}

\subsection{Glosario general}

Las secciones anteriores ya explicaron por separado VLM, CoT, espacio de acciones, aprendizaje en línea y modelo del mundo. Esta sección reúne esos términos para que el lector pueda repasar la cadena argumentativa de todo el episodio: desde las propiedades de los datos de las imágenes estáticas, pasando por la función objetivo de la next token prediction, hasta la reflexión al estilo o1, la Long CoT visual y el aprendizaje en línea.

\noindent\begin{tabularx}{\textwidth}{p{0.18\textwidth}p{0.36\textwidth}X}
\toprule
Término & Definición breve & Papel en este episodio \\
\midrule
VLM & Vision-Language Model, modelo de visión y lenguaje & Forma principal de modelo que conecta imágenes, video y texto. \\
MIM & Masked Image Modeling, modelado de imágenes enmascaradas & Ruta representativa de la CV hacia el aprendizaje por preentrenamiento al estilo BERT. \\
Next Token Prediction & Objetivo de entrenamiento que predice el siguiente token & Explica el origen común de la potencia de los modelos de lenguaje y de su defecto de saltarse pasos. \\
CoT & Chain of Thought, cadena de pensamiento & Despliega una decisión de un solo paso en un proceso en el que se puede buscar. \\
Meta-CoT & Repensar y reflexionar sobre la CoT & Atiende las critical decisions y la elección de rutas. \\
Action Space & Conjunto de acciones que el modelo puede elegir en una tarea & Determina si el RL puede encontrar mediante búsqueda la conducta correcta. \\
Online Learning & Aprendizaje en línea: el sistema se actualiza de forma continua mientras se usa & Podría ser la clave de la siguiente generación de Agents. \\
World Model & Modelo interno que predice cómo cambia el estado del mundo & Conecta lo multimodal, la robótica y la inteligencia corporizada. \\
\bottomrule
\end{tabularx}

\subsection{Preguntas para observar en adelante}

Esta sección no enumera preguntas genéricas: convierte los juicios de la entrevista sobre las rutas técnicas en puntos de observación que se pueden seguir. En los próximos dos años, si de verdad aparece un nuevo momento GPT-4 en lo multimodal, las señales deberían verse en estos puntos: no solo que los modelos suban en los benchmarks, sino que surjan nuevas capacidades de sistema en razonamiento visual, controlabilidad de la generación, long context y aprendizaje en línea.

\begin{enumerate}
\item ¿Puede la Long CoT visual mejorar de forma notable el razonamiento multimodal en benchmarks reales, y no solo generar explicaciones más largas?
\item ¿Dará un salto cualitativo la controlabilidad de los modelos generativos gracias a la CoT visual, la edición local y el entrenamiento con retroalimentación?
\item ¿El long context seguirá dependiendo de ampliar la ventana de un solo modelo, o se moverá hacia esquemas planner/worker, la colaboración entre varios modelos y la consulta mediante herramientas, como quien hojea un libro?
\item ¿Cómo resolverá el aprendizaje en línea los problemas de seguridad, contaminación de datos, retroalimentación escasa y costo de la actualización continua?
\item ¿Podrá la robótica, antes de adelantarse en el hardware y en el cuerpo del robot, completar las capacidades básicas de razonamiento visual, modelo del mundo y retroalimentación de acciones?
\end{enumerate}

\subsection{Lecturas adicionales}

\begin{itemize}
\item Si le interesan el entrenamiento de Agents y la ingeniería de contexto, puede compararlo con las notas del EP110 sobre los informes técnicos de Kimi K2 / ChatGPT Agent / Qwen3-Coder.
\item Si le interesan el Agentic LLM y el escalamiento en tiempo de prueba, puede compararlo con las notas del EP113, la entrevista a Yang Zhilin (杨植麟) sobre Kimi K2.
\item Si le interesan los problemas de datos y de modelos del mundo en la inteligencia corporizada, puede compararlo con las entrevistas del EP106 a Wang He (王鹤), del EP109 a Xie Chen (谢晨) y del EP121 a Tan Jie (谭捷).
\end{itemize}

\end{document}
